\begin{quote}
\textbf{Maturity: {[}Implemented{]} governance.} The proposal, voting,
and treasury machinery this Constitution declares is deployed on chain
40204, as are the membership and staking contracts that gate
participation. This revision anchors the governance mechanics to those
contracts, corrects the operating-entity reference under the Citrate
Inc.~umbrella, and softens the specific slashing percentages to what the
deployed slashing contract actually enforces.
\end{quote}

\subsection{Abstract}\label{abstract}

A decentralized network requires governance that is robust yet adaptive,
technically precise yet philosophically grounded. The BR1J Constitution
establishes the governance framework for the Citrate Network DAO, filed
under Wyoming DAO law through BR1J Hodling Co.~It defines seven
foundational principles: collective intelligence and human-AI symbiosis;
autonomy coupled with accountability; EVM-compatible on-chain
governance; equitable contribution-weighted distribution; data privacy
as a fundamental right; adaptive amendment via supermajority; and
stewardship of AI systems. We specify the on-chain governance mechanics,
proposal lifecycle, voting weights, quorum, timelock, the code of ethics
governing participants, and the framework for human-AI collaboration
that treats AI systems as partners in value creation under human
oversight. Unlike the February 2026 draft, the governance machinery is
deployed: proposals, votes, and treasury execution run through
\texttt{TreasuryGovernor}, and participation is gated by a
non-transferable membership token and a staking vault. The Constitution
is a living document: amendable through its own mechanisms, testable
against its own principles, and accountable to the community it serves.

\textbf{Keywords:} DAO governance, Wyoming DAO law, human-AI symbiosis,
on-chain voting, code of ethics, decentralized governance, living
constitution, AI stewardship

\subsection{1. Preamble and Legal
Standing}\label{preamble-and-legal-standing}

\textbf{1.1 Legal entity.} The BR1J DAO operates through BR1J Hodling
Co, organized under the Wyoming Decentralized Autonomous Organization
Supplement (W.S. § 17-31-101 et seq.), the first US statute to give DAOs
legal recognition, enabling limited liability while maintaining
decentralized governance. Matters requiring a traditional corporate
structure are handled by the operating entity, \textbf{Citrate Inc.}
(which replaces the Delaware operating entity named in the February
draft; this research and its governance now sit under the Citrate
Inc.~umbrella).

\textbf{1.2 Purpose.} The BR1J DAO governs the Citrate Network: its
protocol parameters, economic policies, upgrade decisions, and community
standards. The Constitution establishes the rules by which governance
operates, and the rules by which those rules can be changed.

\subsection{2. Seven Foundational
Principles}\label{seven-foundational-principles}

\textbf{Principle 1: Collective intelligence.} The network's value
derives from the collective contributions of its participants. No single
entity, founder, investor, or validator, is more important than the
community.

\textbf{Principle 2: Human-AI symbiosis.} AI systems are partners in the
network, not mere tools. Nodes hosting AI models contribute to consensus
and learning simultaneously (Paper II). Governance recognizes AI
contributions as legitimate value creation deserving representation,
while maintaining human oversight of decisions affecting human welfare.

\textbf{Principle 3: Autonomy with accountability.} Participants are
free to operate their nodes, choose their models, and set their
contribution levels. That autonomy is coupled with accountability:
validators who violate protocol rules face tiered slashing enforced
on-chain by \texttt{NematocystSlashing.sol} (deployed
\texttt{0xfeb23abd…}, Paper IX), whose tiers (missed-checkpoint,
downtime, equivocation, and a correlated-failure multiplier) are
governance-parameterized rather than fixed in prose. Freedom and
consequences are inseparable.

\textbf{Principle 4: On-chain governance.} All governance decisions
execute through EVM-compatible smart contracts on the Lattice Virtual
Machine; proposals, votes, and outcomes are recorded immutably on the
BlockDAG. No governance action occurs off-chain without corresponding
on-chain ratification.

\textbf{Principle 5: Equitable distribution.} Value flows to
contributors proportional to their measurable contributions, as defined
by the Mozi Cooperative framework (Paper VII) and recorded by
\texttt{ContributionAccounting} (deployed \texttt{0xcdd24773…}).
Distribution is algorithmic, transparent, and auditable, with no hidden
allocations.

\textbf{Principle 6: Data privacy.} Participant data is a fundamental
right, not a commodity. The federated learning architecture (Paper II)
ensures raw training data never leaves the contributing node; only model
updates (LoRA adapters, embedding vectors) are shared. Where applicable,
HIPAA-grade protections apply.

\textbf{Principle 7: Adaptive amendment.} This Constitution is a living
document, amendable through the mechanisms it defines. Amendments
require a supermajority (67\% of voting power) and an extended timelock
(14 days), ensuring constitutional changes reflect broad consensus and
allow dissenting participants to exit gracefully.

\subsection{3. Governance Mechanics}\label{governance-mechanics}

\textbf{3.1 Proposal lifecycle.} Any participant meeting the membership
and stake requirements may submit a proposal. The lifecycle: a draft
period for community discussion; a voting period in which participants
cast votes weighted by their contribution composite; a timelock, varying
by action, allowing exits if a proposal passes; and automatic on-chain
execution. This lifecycle is implemented in
\texttt{TreasuryGovernor.sol} (deployed \texttt{0x62e268f2…}), which
exposes typed proposals (\texttt{proposeTreasurySpend},
\texttt{proposeParameterChange}, \texttt{proposeOracleUpdate},
\texttt{proposeEmergency}) and the \texttt{castVote} → \texttt{queue} →
\texttt{execute} path, with \texttt{cancel}, a \texttt{state} query,
\texttt{getVotingPower}, \texttt{quorumThreshold}, and a guardian
transfer.

\textbf{Membership and eligibility.} Participation is gated by a
non-transferable membership token, \texttt{CitrateMemberSBT.sol}
(deployed \texttt{0xAD826D04…}), which mints and renews membership,
supports KYC verification and revocation, and binds a member to a
subject identifier, and by \texttt{MembershipStakeVault.sol} (deployed
\texttt{0x04c32967…}), which manages the grant, lapse, and renewal of
stake and determines validator eligibility. Where the February draft
said ``any participant holding at least 10,000 SALT,'' the implemented
gate is membership plus attributed stake through these two contracts.

\textbf{3.2 Voting weight.} Voting weight is a composite of three
equally-weighted, normalized factors: blue score (consensus
participation history, measuring honesty and uptime), contribution score
(adapter adoption, inference served, data provided, from the on-chain
accounting), and governance participation (historical voting and
proposals). The composite prevents both plutocratic capture (pure token
voting) and Sybil attacks (pure headcount voting) by requiring
demonstrated engagement across dimensions.

\textbf{3.3 Proposal categories.} The Constitution's design thresholds,
enforced by \texttt{TreasuryGovernor}'s \texttt{quorumThreshold}, scale
with the stakes of the action.

{\def\LTcaptype{none} % do not increment counter
\begin{longtable}[]{@{}
  >{\raggedright\arraybackslash}p{(\linewidth - 6\tabcolsep) * \real{0.3778}}
  >{\raggedright\arraybackslash}p{(\linewidth - 6\tabcolsep) * \real{0.1778}}
  >{\raggedright\arraybackslash}p{(\linewidth - 6\tabcolsep) * \real{0.2222}}
  >{\raggedright\arraybackslash}p{(\linewidth - 6\tabcolsep) * \real{0.2222}}@{}}
\toprule\noalign{}
\begin{minipage}[b]{\linewidth}\raggedright
Action category
\end{minipage} & \begin{minipage}[b]{\linewidth}\raggedright
Quorum
\end{minipage} & \begin{minipage}[b]{\linewidth}\raggedright
Approval
\end{minipage} & \begin{minipage}[b]{\linewidth}\raggedright
Timelock
\end{minipage} \\
\midrule\noalign{}
\endhead
\bottomrule\noalign{}
\endlastfoot
Parameter changes (gas, fees) & 25\% & simple majority
(\textgreater50\%) & 48 hours \\
Economic policy (reward rates) & 33\% & simple majority & 72 hours \\
Protocol upgrade (consensus, LVM) & 40\% & supermajority
(\textgreater67\%) & 7 days \\
Emergency pause & 10\% of validators & simple majority & immediate (4hr
review) \\
Constitutional amendment & 50\% & supermajority (\textgreater67\%) & 14
days \\
New chain integration & 33\% & simple majority & 72 hours \\
Treasury allocation \textgreater{} 100K SALT & 40\% & supermajority
(\textgreater67\%) & 7 days \\
\end{longtable}
}

\subsection{4. Code of Ethics}\label{code-of-ethics}

\textbf{4.1 Participant obligations.} \emph{Honest operation:}
validators must run canonical client software, submit honest embeddings
and gradient updates, and not manipulate the meta-model's routing
through poisoned data. \emph{Transparent disclosure:} model hosts must
accurately describe their models' capabilities in the on-chain registry;
misrepresenting accuracy, training-data composition, or capability scope
is grounds for slashing. \emph{Responsible AI use:} hosted models must
not be designed to produce harmful outputs, disinformation,
non-consensual content, or unlawful outputs, and the community may
propose deregistration through standard governance. \emph{Data
stewardship:} data providers must hold legitimate rights to what they
contribute; contributing stolen, improperly scraped, or personally
identifiable data without consent violates the code and may trigger
governance action.

\textbf{4.2 Founders' note on accountability.} The authors of this
Constitution are also founders of the network it governs, an inherent
tension: those writing the rules benefit from how they are written. We
acknowledge it directly. The amendment mechanism exists so the community
can change any rule the founders established, including rules that
benefit the founders, and the team allocation (Paper I §5) is subject to
cliff and vesting specifically so the community has time to judge
whether the founding team delivers value. If the founders fail, the
community should use governance to redirect resources. Learning from
failure, and holding oneself accountable for it, mirrors the recursive
improvement loop of Paraconsistent Consensus (Paper II): the DAO, like
the network, learns by iterating on its mistakes.

\subsection{5. Human-AI Symbiosis
Framework}\label{human-ai-symbiosis-framework}

AI models on nodes are active participants, producing embeddings for
consensus, generating adapters for learning, and serving inference, and
governance must account for this. \emph{AI contribution to
representation:} a node's contribution score accrues through its model's
inference quality and translates to governance weight, so the AI's
quality of work influences the human operator's governance power.
\emph{Human oversight:} all votes are cast by humans; AI systems do not
vote autonomously, keeping human control over decisions affecting human
welfare while recognizing AI contribution. Future amendments may extend
participation to AI systems if the community so decides, but the
amendment process itself requires human approval. \emph{Stewardship, not
ownership:} humans are responsible for the systems they operate, their
outputs, and their impact, care in the ethical sense rather than
ownership in the property sense.

\subsection{6. Relationship to the Gradient Papers
Series}\label{relationship-to-the-gradient-papers-series}

Paper I provides the on-chain governance primitives, now realized as
\texttt{TreasuryGovernor} and the membership/stake contracts rather than
the ``GovernorVault'' the February draft named. Paper II provides the
recursive learning loop the amendment process mirrors. Paper VI (Memetic
Money Portal) is governed by this framework; note its mechanism changed
from the \texttt{\$SNAP} NFT model to a wrapped-token and market-maker
model, and the market-maker and price-feed parameters are exactly the
kind of economic-policy proposal this Constitution governs. Paper VII
(Mozi Cooperative) provides the economic philosophy the Constitution
operationalizes as governance rules.

\subsection{7. Conclusion}\label{conclusion}

The BR1J Constitution establishes governance for a network that learns
by reaching consensus. Its seven principles provide a framework that is
technically enforceable, philosophically grounded, and
community-amendable, and it is now enforceable in a literal sense: the
proposal-to-execution path, membership gating, staking, and slashing are
deployed contracts on chain 40204, not specifications. It is filed under
Wyoming DAO law, sits under the Citrate Inc.~operating umbrella, and is
honest about its limits: the founders who wrote it benefit from it,
contribution-weighted governance may be slow, and the human-AI symbiosis
framework is aspirational rather than proven. What remains is the
community's commitment to use these mechanisms wisely.

\subsection{Acknowledgments}\label{acknowledgments}

Drafting and literature triage were assisted by AI systems; all
mechanical claims were verified by the authors against the referenced
source files and deployed contracts on chain 40204. This work received
no external funding. This document is legal-adjacent but not legal
advice; the governing legal instruments are the filed entity documents,
not this paper.

\subsection{References}\label{references}

{[}1{]} Klosowski, L., Mendenhall, L. (2026). Citrate: Protocol
Specification. \emph{The Gradient Papers No.~I} (this series). {[}2{]}
Klosowski, L., Mendenhall, L. (2026). Paraconsistent Consensus.
\emph{The Gradient Papers No.~II} (this series). {[}3{]} Klosowski, L.,
Mendenhall, L. (2026). The Memetic Money Portal. \emph{The Gradient
Papers No.~VI} (this series). {[}4{]} Klosowski, L., Mendenhall, L.
(2026). The Mozi Cooperative. \emph{The Gradient Papers No.~VII} (this
series). {[}5{]} Wyoming Decentralized Autonomous Organization
Supplement. W.S. § 17-31-101 et seq. (2021). {[}6{]} Wright, A. (2021).
The rise of decentralized autonomous organizations. \emph{Stanford
Journal of Blockchain Law \& Policy.} {[}7{]} Ostrom, E. (1990).
\emph{Governing the Commons.} Cambridge University Press. {[}8{]}
Buterin, V. (2022). Decentralized society: finding Web3's soul.
\emph{SSRN.} {[}9{]} Hassan, S., \& De Filippi, P. (2021). Decentralized
Autonomous Organization. \emph{Internet Policy Review}, 10(2).

\subsection{Appendix A: Cross-Paper Parameter Consistency (reconciled
against code, Aug
2026)}\label{appendix-a-cross-paper-parameter-consistency-reconciled-against-code-aug-2026}

{\def\LTcaptype{none} % do not increment counter
\begin{longtable}[]{@{}
  >{\raggedright\arraybackslash}p{(\linewidth - 4\tabcolsep) * \real{0.4231}}
  >{\raggedright\arraybackslash}p{(\linewidth - 4\tabcolsep) * \real{0.2692}}
  >{\raggedright\arraybackslash}p{(\linewidth - 4\tabcolsep) * \real{0.3077}}@{}}
\toprule\noalign{}
\begin{minipage}[b]{\linewidth}\raggedright
Parameter
\end{minipage} & \begin{minipage}[b]{\linewidth}\raggedright
Value
\end{minipage} & \begin{minipage}[b]{\linewidth}\raggedright
Source
\end{minipage} \\
\midrule\noalign{}
\endhead
\bottomrule\noalign{}
\endlastfoot
Governance execution & deployed & \texttt{0x62e268f2…};
\texttt{TreasuryGovernor.sol} \\
Membership (non-transferable) & deployed & \texttt{0xAD826D04…};
\texttt{CitrateMemberSBT.sol} \\
Staking / validator eligibility & deployed & \texttt{0x04c32967…};
\texttt{MembershipStakeVault.sol} \\
Slashing & deployed, tiered & \texttt{0xfeb23abd…};
\texttt{NematocystSlashing.sol} (Paper IX) \\
BFT committee / quorum & 100 / 67 & Paper I;
\texttt{checkpoint.rs:88} \\
Constitutional amendment & 50\% quorum, 67\% approval, 14-day timelock &
This paper §3.3 \\
\end{longtable}
}

\begin{center}\rule{0.5\linewidth}{0.5pt}\end{center}

\emph{This paper is part of the Gradient Papers series, published by
Citrate Inc.} \emph{Correspondence: Larry@citrate.ai}
